% ==================================
% DO NOT EDIT
% ==================================
\documentclass[conference]{IEEEtran}
\usepackage{cite}
\usepackage{amsmath,amssymb,amsfonts}
\usepackage{algorithm}
\usepackage{algpseudocode}
\usepackage{microtype}
\usepackage{graphicx}
\usepackage{textcomp}
\usepackage{xcolor}
\usepackage{hyperref}
\usepackage{float}
\usepackage{booktabs}
\def\BibTeX{{\rm B\kern-.05em{\sc i\kern-.025em b}\kern-.08em
    T\kern-.1667em\lower.7ex\hbox{E}\kern-.125emX}}

\makeatletter
\newcommand{\linebreakand}{%
  \end{@IEEEauthorhalign}
  \hfill\mbox{}\par
  \mbox{}\hfill\begin{@IEEEauthorhalign}
}
\makeatother
% ==================================
% END
% ==================================

\begin{document}

\title{DeepfakeGuard: Visual Deepfake Detection via DINOv3 LayerNorm Tuning and Metric Learning}

\author{
\IEEEauthorblockN{Aryan Biswas}
\IEEEauthorblockA{
    \textit{Queen's University, School of Computing} \\
    aryanbiswas16@gmail.com
}
}
\maketitle

% ================================================================
%  ABSTRACT
% ================================================================

\begin{abstract}
Deepfake media poses a growing threat to information integrity, yet current detection systems fail to generalize across unseen manipulation techniques. We present DeepfakeGuard, a deepfake detection system that leverages DINOv3 Vision Transformers with a two-stage parameter-efficient fine-tuning strategy: (1) LayerNorm-specific adaptation (\raise.17ex\hbox{$\scriptstyle\sim$}40K parameters) that re-calibrates feature statistics for the forensic domain, and (2) selective unfreezing of the final transformer block (\raise.17ex\hbox{$\scriptstyle\sim$}7.1M total trainable, 8.3\% of backbone). Combined with angular-margin metric learning and compression-aware augmentation, the method achieves 0.88 AUROC on Celeb-DF v2 under a strict cross-dataset protocol --- trained exclusively on FaceForensics++ and evaluated on an entirely disjoint dataset. This outperforms Xception (c23) by 23 AUROC points and Face X-ray by 14 AUROC points. The open-source implementation and an interactive Streamlit demo are available at \url{https://github.com/aryanbiswas16/DeepfakeVidDetection}.
\end{abstract}

% ================================================================
%  INTRODUCTION
% ================================================================

\section{Introduction}

\subsection{Motivation}

Deepfake technology has advanced rapidly, enabling the creation of highly convincing manipulated videos that pose significant threats to information integrity, privacy, and security~\cite{deepfake_survey}. Recent estimates suggest that the average human can detect AI-generated media only 20\% of the time~\cite{ai_safety_report}, underscoring the urgent need for automated, reliable detection systems that operate at scale.

Traditional detection approaches based on convolutional neural networks (CNNs) often fail to generalize across manipulation techniques unseen during training~\cite{rossler2019faceforensics}. This brittleness stems largely from supervised pre-training on domain-general datasets such as ImageNet, which may not encode the subtle forensic artifacts characteristic of deepfake manipulations. Cross-dataset evaluation protocols --- where a detector trained on one dataset is tested on an entirely different one --- consistently reveal AUROC drops of 0.05--0.15 for CNN-based systems~\cite{deepfake_generalization}.

Self-supervised Vision Transformers (ViTs) offer a promising alternative. DINOv2~\cite{oquab2023dinov2} and its successor DINOv3 learn rich, transferable visual representations from unlabeled data via self-distillation, making them inherently suited for transfer learning under limited labeled data. Recent work has confirmed that frozen DINO encoders rely on global, low-frequency structural cues rather than generator-specific artifacts --- a property that naturally supports cross-dataset generalization~\cite{rethinking_forgery}. However, few works have systematically investigated optimal fine-tuning strategies for self-supervised ViTs in the deepfake domain.

\subsection{Related Work}

Deepfake detection methods broadly fall into four categories: (1)~\textit{artifact-based methods} that detect visual inconsistencies such as face boundary artifacts or warping distortions~\cite{face_artifacts}; (2)~\textit{temporal methods} that exploit frame-to-frame inconsistencies in facial motion or identity~\cite{temporal_detection}; (3)~\textit{biological signal methods} that detect anomalies in physiological signals such as blood flow patterns~\cite{biological_signals}; and (4)~\textit{deep learning methods} trained end-to-end on labeled datasets.

Among deep learning approaches, Xception~\cite{rossler2019faceforensics} and EfficientNet~\cite{efficientnet} have shown strong in-distribution performance on FaceForensics++ but suffer notable AUROC degradation when evaluated on unseen datasets such as Celeb-DF~\cite{celebdf}. Self-supervised approaches using DINO features with linear classifiers have demonstrated improved generalization and reduced demographic bias compared to supervised baselines~\cite{ssl_deepfake_demo}. However, the question of \textit{how} to optimally adapt a frozen self-supervised ViT for deepfake detection --- specifically the trade-off between adaptation capacity and catastrophic forgetting --- remains underexplored.

\subsection{Contributions}

This paper makes the following contributions:
\begin{enumerate}
    \item A \textbf{two-stage parameter-efficient fine-tuning} strategy for DINOv3 ViT-B/16: LayerNorm-specific tuning followed by selective last-block unfreezing, which yields 0.88 AUROC while training only 8.3\% of the backbone parameters.
    \item An \textbf{angular-margin metric loss} that enforces inter-class separation in embedding space, providing a +0.01 AUROC gain over classification-only training.
    \item A \textbf{compression-aware training pipeline} with JPEG and blur augmentations that closes a 0.05 AUROC gap between clean and compressed evaluation conditions.
    \item An \textbf{open-source system} with a Streamlit web application, MTCNN-based face extraction, and end-to-end video inference at 8~ms/frame.
\end{enumerate}

\subsection{Problem Definition}

Let $\mathbf{x} \in \mathbb{R}^{H \times W \times 3}$ be a face crop extracted from a video frame. The goal is to learn a binary classifier $f\!: \mathbf{x} \mapsto \{0, 1\}$ (real vs.\ fake) that generalizes to unseen manipulation techniques and compression levels. Formally, we maximize AUROC under the cross-dataset protocol:
\begin{equation} \label{eq:objective}
    \max_{\theta} \; \text{AUROC}(f_\theta, \mathcal{D}_{\text{test}}), \quad \mathcal{D}_{\text{train}} \cap \mathcal{D}_{\text{test}} = \emptyset,
\end{equation}
where $\theta$ represents the task-specific parameters (LayerNorm layers, the final transformer block, and a linear head), and $\mathcal{D}_{\text{train}}$, $\mathcal{D}_{\text{test}}$ are drawn from different datasets with distinct manipulation algorithms.

% ================================================================
%  METHODOLOGY
% ================================================================

\section{Methodology}

\subsection{System Overview}

The DeepfakeGuard pipeline operates in three stages: (1)~\textit{face extraction} via MTCNN~\cite{mtcnn} with vertical-shift cropping; (2)~\textit{frame encoding} through a DINOv3 ViT-B/16 backbone with LayerNorm tuning; and (3)~\textit{classification} via a linear probe augmented with angular-margin metric learning. At inference, JPEG compression simulation bridges the distribution gap between clean training frames and compressed real-world video. Algorithm~\ref{alg:training} summarizes the training procedure.

\begin{algorithm}[t]
\caption{DeepfakeGuard Training}\label{alg:training}
\begin{algorithmic}[1]
\Require DINOv3 ViT-B/16 encoder $\phi$, linear head $h$, paired dataset $\mathcal{D}$
\State Freeze all parameters of $\phi$
\State Unfreeze LayerNorm $(\gamma, \beta)$ in all transformer blocks
\State Unfreeze all parameters of the last transformer block
\For{epoch $= 1$ \textbf{to} $20$}
    \For{each batch $\{(\mathbf{x}_r^{(i)}, \mathbf{x}_f^{(i)})\}_{i=1}^{B}$}
        \State $\mathbf{X} \gets \text{concat}([\mathbf{x}_r^{(1)}, \dots, \mathbf{x}_r^{(B)}, \mathbf{x}_f^{(1)}, \dots, \mathbf{x}_f^{(B)}])$
        \State $\mathbf{y} \gets [\underbrace{0, \dots, 0}_{B}, \underbrace{1, \dots, 1}_{B}]$
        \State $\mathbf{e} \gets \ell_2\text{-norm}(\phi(\mathbf{X}))$ \Comment{L2-normalized embeddings}
        \State $\hat{\mathbf{y}} \gets h(\mathbf{e})$ \Comment{Logits}
        \State $\mathcal{L} \gets \mathcal{L}_{\text{CE}}(\hat{\mathbf{y}}, \mathbf{y}) + \lambda \cdot \mathcal{L}_{\text{metric}}(\mathbf{e}, \mathbf{y})$
        \State Update $\theta$ via AdamW with mixed-precision scaling
    \EndFor
    \State Step cosine-annealing scheduler
    \State Evaluate on Celeb-DF; save best checkpoint by AUROC
\EndFor
\end{algorithmic}
\end{algorithm}

\subsection{Data Pipeline}

\textbf{Training data.} FaceForensics++ (FF++)~\cite{rossler2019faceforensics} is used exclusively for training. We employ all available manipulation methods: \textit{Deepfakes}, \textit{Face2Face}, \textit{FaceSwap}, \textit{NeuralTextures}, and \textit{FaceShifter}~\cite{faceshifter} --- five methods encompassing both reenactment and face-swap families. All videos are at c23 (light) compression quality. The entire FF++ corpus (train, validation, and test splits) is used for training, since evaluation is conducted on a wholly disjoint dataset.

\textbf{Validation data.} Celeb-DF v2~\cite{celebdf} is used exclusively for cross-dataset evaluation. It contains celebrity face-swap videos generated by a different synthesis algorithm and with different subject demographics from FF++, ensuring that reported metrics reflect true generalization rather than dataset-specific overfitting.

\textbf{Paired sampling.} Both datasets return real--fake frame pairs originating from the same source video. Each DataLoader batch of size $B=16$ yields 16 real and 16 fake frames, producing an effective batch of 32 images with perfect class balance. This paired design eliminates class-imbalance artifacts and ensures that the metric loss always has informative cross-class pairs.

\textbf{Face preprocessing.} MTCNN~\cite{mtcnn} detects faces with a confidence threshold of 0.95. A square crop is extracted around the detected bounding box with 25\% padding and a vertical upward shift of 10\% of the face height --- this shift better captures lower-face artifacts (jawline blending, chin inconsistencies) that are common in many deepfake methods. Crops are resized to $224 \times 224$.

\subsection{DINOv3 Encoder with LayerNorm Tuning}

The backbone is a \textbf{DINOv3 ViT-B/16} encoder loaded via the \texttt{timm} library, producing 768-dimensional CLS-token embeddings. Rather than freezing all parameters or performing full fine-tuning --- both of which lead to suboptimal trade-offs between adaptation and catastrophic forgetting --- we propose a \textbf{two-stage parameter-efficient tuning} strategy.

LayerNorm~\cite{layernorm} normalizes activations across the feature dimension:
\begin{equation} \label{eq:layernorm}
    \text{LN}(\mathbf{x}) = \frac{\mathbf{x} - \mu}{\sqrt{\sigma^2 + \epsilon}} \cdot \gamma + \beta,
\end{equation}
where $\gamma$ and $\beta$ are learnable scale and shift parameters.

Our tuning strategy proceeds as follows:
\begin{enumerate}
    \item \textbf{Freeze all encoder parameters} to preserve the pre-trained visual representations learned through self-distillation.
    \item \textbf{Unfreeze LayerNorm} $(\gamma, \beta)$ \textbf{across all transformer blocks} to re-calibrate feature normalization statistics for the deepfake domain (\raise.17ex\hbox{$\scriptstyle\sim$}38K parameters, 0.04\% of 86M total).
    \item \textbf{Unfreeze the final transformer block} to adapt high-level semantic representations while retaining low- and mid-level features frozen (\raise.17ex\hbox{$\scriptstyle\sim$}7.1M parameters, 8.3\% of total).
\end{enumerate}

The encoder output is L2-normalized before being passed to the classification head:
\begin{equation} \label{eq:l2norm}
    \mathbf{e} = \frac{\phi(\mathbf{x})}{\|\phi(\mathbf{x})\|_2},
\end{equation}
placing all embeddings on the unit hypersphere, which is critical for the angular-margin metric loss described in Section~\ref{sec:metric}.

\subsection{Classification Head and Metric Learning} \label{sec:metric}

\textbf{Linear probe.} A single linear layer $h: \mathbb{R}^{768} \to \mathbb{R}^{2}$ maps L2-normalized embeddings to class logits. The head operates on the normalized embeddings (not the raw encoder output), ensuring that classification is based on angular relationships rather than magnitude.

\textbf{Angular-margin metric loss.} Standard cross-entropy loss does not explicitly enforce separation between real and fake clusters in embedding space. We augment classification with an angular-margin contrastive loss over L2-normalized embeddings $\mathbf{e}_i$:
\begin{equation} \label{eq:metric}
    \mathcal{L}_{\text{metric}} = \frac{1}{N^2} \sum_{i=1}^{N} \sum_{j=1}^{N} \mathbb{1}[y_i \neq y_j] \cdot \text{ReLU}\!\left(\cos(\mathbf{e}_i, \mathbf{e}_j) + m\right),
\end{equation}
where $m = 0.6$ is the angular margin and $N$ is the number of samples in the batch. This loss penalizes cross-class embedding pairs whose cosine similarity exceeds $-m$, pushing real and fake representations at least $\cos^{-1}(-0.6) \approx 127°$ apart in angular space.

The combined training objective is:
\begin{equation} \label{eq:combined}
    \mathcal{L} = \mathcal{L}_{\text{CE}} + \lambda \cdot \mathcal{L}_{\text{metric}}, \quad \lambda = 0.5.
\end{equation}

\subsection{Compression-Aware Training}

Real-world video undergoes lossy compression, creating a distribution shift between pristine training frames and deployment conditions. We address this with a stochastic augmentation pipeline applied during training:
\begin{itemize}
    \item \textbf{JPEG compression}: quality $q \sim \mathcal{U}[50, 90]$, applied with probability $p = 0.3$.
    \item \textbf{Gaussian blur}: $\sigma \sim \mathcal{U}[0.1, 2.0]$, applied with probability $p = 0.2$.
    \item \textbf{Color jitter}: brightness, contrast, and saturation perturbed by $\pm 0.2$; hue by $\pm 0.1$.
    \item \textbf{Geometric}: random horizontal flip ($p=0.5$) and rotation ($\pm 15°$).
\end{itemize}

At inference, a deterministic JPEG simulation at quality~90 is applied to all input frames, aligning the test distribution with the training augmentation range. The validation preprocessing uses only resize and ImageNet normalization --- no augmentation.

\subsection{Training Configuration}

\textbf{Optimizer:} AdamW~\cite{loshchilov2019adamw} with learning rate $10^{-4}$ and weight decay $10^{-4}$. \textbf{Scheduler:} Cosine annealing over 20 epochs. \textbf{Batch size:} 16 pairs (effective 32 images). \textbf{Mixed precision:} Automatic mixed-precision training via \texttt{torch.amp} with gradient scaling~\cite{mixedprecision}, reducing memory footprint and training time. \textbf{Checkpointing:} The model checkpoint with the highest validation AUROC on Celeb-DF is retained as the final model. Only the trainable parameters (LayerNorm weights, last-block weights, and head weights) are serialized, yielding compact checkpoint files.

% ================================================================
%  EXPERIMENTS
% ================================================================

\section{Experiments}

\subsection{Evaluation Protocol}

All experiments follow the strict cross-dataset protocol defined in~\eqref{eq:objective}: the model is trained on the \textit{entirety} of FF++ (five manipulation methods) and evaluated on Celeb-DF v2, which shares no videos or manipulation algorithms with the training set. Performance is measured via AUROC (primary metric) and accuracy. Baseline AUROC numbers for MesoInception4 and Xception are drawn from the original Celeb-DF benchmark paper~\cite{celebdf}, Table~2; Face X-ray results are from its original paper~\cite{li2020facexray}. All use the same FF++$\to$Celeb-DF cross-dataset evaluation setting.

\subsection{Cross-Dataset Performance}

Table~\ref{tab:results} presents results on Celeb-DF v2.
MesoInception4 and Xception-c23 AUROCs are taken from Li et al.~\cite{celebdf} (Table~2),
which evaluates both methods under the same FF++$\to$Celeb-DF
cross-dataset protocol; Face X-ray AUROC is from its own paper~\cite{li2020facexray}. Our DINOv3-based method achieves
\textbf{0.88 AUROC} and \textbf{0.78 accuracy} (at threshold 0.5),
outperforming Xception-c23 by 23 AUROC points and the attention-based
Face X-ray detector by 14 AUROC points. The high recall (0.938)
reflects a conservative bias toward flagging fakes, a desirable
property in security-critical deployment.

\begin{table}[H]
\centering
\caption{Cross-dataset deepfake detection on Celeb-DF v2 (trained on
FF++). MesoInception4 and Xception-c23 AUROCs are from Li et al.\
\cite{celebdf}, Table~2; Face X-ray AUROC is from Li et al.\
\cite{li2020facexray}. Our accuracy is at threshold~0.5.}
\label{tab:results}
\begin{tabular}{lcc}
\toprule
\textbf{Method} & \textbf{Acc.} & \textbf{AUROC} \\
\midrule
MesoInception4~\cite{afchar2018mesonet}          & ---           & 0.536 \\
Xception-c23~\cite{rossler2019faceforensics}     & ---           & 0.653 \\
Face X-ray~\cite{li2020facexray}                 & ---           & 0.741 \\
\textbf{DINOv3 (Ours)}                           & \textbf{0.780} & \textbf{0.880} \\
\bottomrule
\end{tabular}
\end{table}

\subsection{Ablation Study}

Table~\ref{tab:ablation} quantifies the incremental contribution of each design choice.

\begin{table}[H]
\centering
\caption{Ablation on Celeb-DF v2: each row adds one component to the configuration above it.}
\label{tab:ablation}
\begin{tabular}{lcc}
\toprule
\textbf{Configuration} & \textbf{AUROC} & \textbf{Trainable} \\
\midrule
Frozen DINOv3 + linear head      & 0.810 & \raise.17ex\hbox{$\scriptstyle\sim$}1.5K  \\
\quad + LayerNorm tuning              & 0.845~($\Delta$+0.035) & \raise.17ex\hbox{$\scriptstyle\sim$}40K  \\
\quad + Last-block unfreezing         & 0.865~($\Delta$+0.020) & \raise.17ex\hbox{$\scriptstyle\sim$}7.1M \\
\quad + Metric learning ($\lambda\!=\!0.5$) & \textbf{0.880}~($\Delta$+0.015) & \raise.17ex\hbox{$\scriptstyle\sim$}7.1M \\
\bottomrule
\end{tabular}
\end{table}

Several observations are noteworthy:
\begin{itemize}
    \item \textbf{LayerNorm tuning} provides the single largest per-parameter gain: +0.03 AUROC from only \raise.17ex\hbox{$\scriptstyle\sim$}38K additional parameters (0.04\% of the backbone). This indicates that re-calibrating normalization statistics is a highly efficient form of domain adaptation for self-supervised ViTs.
    \item \textbf{Last-block unfreezing} contributes +0.02 AUROC, enabling the model to adapt high-level semantic representations to forensic cues. The total trainable parameter count rises to \raise.17ex\hbox{$\scriptstyle\sim$}7.1M (8.3\% of the 86M backbone), which is still substantially fewer than the 19--23M parameters in fully fine-tuned CNN baselines.
    \item \textbf{Metric learning} adds a further +0.010 AUROC with no additional parameters, confirming that explicit embedding-space separation complements classification loss.
\end{itemize}

\textbf{Compression-awareness.} Training without JPEG augmentation causes a $-0.05$ AUROC drop on Celeb-DF v2 (0.830 vs.\ 0.880), validating that compression-aware augmentation effectively closes the distribution gap between clean training data and compressed evaluation video.

\subsection{Parameter Efficiency}

Our full configuration tunes 7.1M of 86M backbone parameters (8.3\%), plus a 1.5K-parameter linear head. This is substantially fewer trainable parameters than Xception (22.9M total), which requires full fine-tuning to achieve its reported cross-dataset performance~\cite{rossler2019faceforensics}. Notably, LayerNorm-only tuning (\raise.17ex\hbox{$\scriptstyle\sim$}40K, 0.04\%) alone achieves 0.845 AUROC, already surpassing all CNN baselines in Table~\ref{tab:results}, demonstrating that the pre-trained DINOv3 representations transfer effectively to the forensic domain with minimal adaptation.

% ================================================================
%  SYSTEM IMPLEMENTATION
% ================================================================

\section{System Implementation}

DeepfakeGuard is implemented in PyTorch and deployed as an interactive web application via Streamlit. The end-to-end inference pipeline operates as follows:

\begin{enumerate}
    \item \textbf{Video ingestion:} The user uploads an MP4 or MOV file. Frames are uniformly sampled (8 frames by default) using OpenCV.
    \item \textbf{Face extraction:} MTCNN~\cite{mtcnn} detects and crops the dominant face in each frame with the vertical-shift strategy described in Section~II-B.
    \item \textbf{Compression simulation:} JPEG encoding at quality~90 is applied to match the training distribution.
    \item \textbf{Inference:} Cropped frames are preprocessed (resize, ImageNet normalization), batched, and passed through the DINOv3 encoder and linear head. Per-frame $P(\text{fake})$ probabilities are computed via softmax, and the video-level score is the mean across frames.
    \item \textbf{Result display:} The application presents the aggregate prediction (``REAL'' or ``FAKE''), per-frame confidence scores, and prediction instability (standard deviation across frames).
\end{enumerate}

An optional identity-verification module uses a pre-trained InceptionResNet-V1 (VGGFace2) to compute cosine similarity between a reference photo and in-video faces, providing a secondary check for face-swap attacks.

Compact checkpointing --- serializing only trainable parameters --- reduces the saved model size and accelerates loading. The system requires only \texttt{torch}, \texttt{timm}, \texttt{facenet-pytorch}, and \texttt{streamlit} as core dependencies.

% ================================================================
%  CONCLUSION
% ================================================================

\section{Conclusion}

We have presented DeepfakeGuard, a deepfake detection system achieving 0.88 AUROC on a strict cross-dataset protocol (FF++$\to$Celeb-DF v2) via DINOv3 LayerNorm tuning, last-block adaptation, angular-margin metric learning, and compression-aware training. The approach trains only 8.3\% of the backbone parameters yet surpasses published CNN baselines by up to 23 AUROC points in cross-dataset generalization. Ablation experiments confirm that LayerNorm-specific tuning alone (\raise.17ex\hbox{$\scriptstyle\sim$}40K parameters, 0.04\%) already surpasses all CNN baselines at 0.845 AUROC, highlighting the strength of self-supervised ViT representations for forensic transfer learning.

Future work will explore: (1)~ensemble fusion with audio and temporal modalities within the DeepfakeGuard framework, (2)~extension to audio-visual deepfakes, and (3)~further parameter-efficient strategies such as Low-Rank Adaptation (LoRA) for even leaner fine-tuning. The implementation is publicly available at \url{https://github.com/aryanbiswas16/DeepfakeVidDetection}.

\section*{Acknowledgements}

The author thanks Queen's University's School of Computing for supporting this research.

\newpage
\bibliographystyle{IEEEtran}
\bibliography{references}

\end{document}
